\subsection{Language Features}

Narrative texts commonly use several grammatical features.

\subsubsection{Simple Past Tense}

The simple past tense is frequently used because narratives commonly describe events that occurred in the past.

\begin{quotation}
  The tortoise walked slowly toward the finish line.
\end{quotation}

\begin{quotation}
  The hare stopped under a tree.
\end{quotation}

Common forms include:

\[
\text{V2}
\]

and

\[
\text{was/were + complement}
\]

\subsubsection{Action Verbs}

Action verbs describe events and activities performed by characters.


Examples include:

\begin{quotation}
  ran, walked, opened, shouted, climbed, escaped, attacked
\end{quotation}

\subsubsection{Temporal Connectives}

Temporal connectives indicate the order or relationship between events.


Examples include:

\begin{quotation}
  first, then, next, after that, suddenly, finally, eventually
\end{quotation}


For example:

\begin{quotation}
  First, the hare ran ahead. Then, he stopped to rest. Finally, the tortoise reached the finish line.
\end{quotation}

\subsubsection{Specific Participants}

Narratives commonly refer to specific characters rather than general groups.


Examples include:

\begin{quotation}
  Belle, the Beast, the tortoise, the hare, a king, a farmer
\end{quotation}

\subsubsection{Descriptive Language}

Adjectives and adverbs are often used to describe characters, settings, and actions.


For example:

\begin{quotation}
  The frightened girl entered the dark and mysterious castle.
\end{quotation}
